% Full CV: Aakarsh Anand (complete record, no page limit). The one-page resume is resume.tex.
% Build: make -C cv

\documentclass[10pt,letterpaper]{article}
\usepackage[margin=0.75in,top=0.55in,bottom=0.62in]{geometry}
% Shared look and header for resume.tex and cv.tex. Change styling or contact details here only.
% EB Garamond for reading text; Gill Sans (a macOS system font) for the name, headings, dates, and metadata.
% No italics anywhere; venues are bold.

\usepackage[lining,scaled=1.06]{ebgaramond}
\setsansfont{Gill Sans}
\usepackage{microtype}
\usepackage{xcolor}
\usepackage{enumitem}
\usepackage{titlesec}
\usepackage{textcase}
\usepackage{hyperref}

\definecolor{accent}{HTML}{3A6532}
\definecolor{muted}{HTML}{5F665F}
\definecolor{rulecolor}{HTML}{BCCFB5}

\hypersetup{colorlinks=true, urlcolor=accent, linkcolor=accent, pdfauthor={Aakarsh Anand}}

\pagestyle{empty}
\hyphenation{RECOMB NeurIPS}
\setlength{\parindent}{0pt}
\setlength{\parskip}{0pt}

% Section headings: small letterspaced sans capitals over a hairline
\newcommand{\secnote}{}
\newcommand{\sectionnote}[1]{\gdef\secnote{#1}}
\newcommand{\headwithrule}[1]{{\addfontfeature{LetterSpace=12}\MakeTextUppercase{#1}}%
  \hspace{0.7em}{\color{rulecolor}\leaders\hrule height 0.72ex depth -0.6ex\hfill}%
  \ifx\secnote\empty\else\hspace{0.7em}{\normalfont\sffamily\footnotesize\color{muted}\secnote}\global\let\secnote\empty\fi}
\titleformat{\section}[block]{\sffamily\color{accent}}{}{0pt}{\headwithrule}
\titlespacing*{\section}{0pt}{8pt}{3.5pt}

\setlist[itemize]{leftmargin=1em, labelsep=0.5em, itemsep=1.5pt, parsep=0pt, topsep=2pt, partopsep=0pt,
  label={\raisebox{0.15ex}{\scriptsize\color{muted}\textbullet}}}

% ---------- macros ----------
\newcommand{\me}{\textcolor{black}{\textbf{A.~Anand}}}
\newcommand{\dates}[1]{{\sffamily\small\color{muted}#1}}
\newcommand{\note}[1]{{\color{muted}#1}}
\newcommand{\vn}[1]{{\color{accent}\textbf{#1}}}
% Separator dot: binds to the preceding word, so a line can end with it but never start with it
\newcommand{\sep}{\unskip\nobreak\enspace{\color{muted}\textperiodcentered}\hspace{0.5em}\ignorespaces}
% \role{title}{organization}{dates}
\newcommand{\role}[3]{\textbf{#1}, #2\hfill\dates{#3}\par}
% \thread{name}{optional note}
\newcommand{\thread}[2]{\vspace{3pt}{\sffamily\color{accent}#1}%
  \if\relax\detokenize{#2}\relax\else{\sffamily\small\color{muted}\sep #2}\fi\par}
% \pub{title}{authors}{venue}{url}: title on one line, authors and venue beneath
\newcommand{\pub}[4]{\item \href{#4}{\textcolor{black}{#1}}\\{\small\color{muted}#2}\sep{\small #3}}

% ---------- header ----------
\newcommand{\header}{%
  \begin{minipage}[b]{0.5\textwidth}
    {\sffamily\fontsize{23}{25}\selectfont\color{accent}\addfontfeature{LetterSpace=9}AAKARSH ANAND}
  \end{minipage}%
  \begin{minipage}[b]{0.5\textwidth}
    \raggedleft\sffamily\small
    \href{mailto:aakarsh@cs.ucla.edu}{aakarsh@cs.ucla.edu}\sep
    \href{https://aakarsh-anand.github.io}{aakarsh-anand.github.io}\\[1pt]
    \href{https://scholar.google.com/citations?user=icSvTl8AAAAJ}{Google Scholar}\sep
    \href{https://github.com/aakarsh-anand}{GitHub}\sep
    \href{https://www.linkedin.com/in/aakarsh-anand-ml/}{LinkedIn}
  \end{minipage}\par
  \vspace{4pt}
  {\color{muted}Machine learning researcher working on foundation models and statistical genetics for
  health data}\par
  \vspace{-2pt}}

\hypersetup{pdftitle={Aakarsh Anand - CV}}

\usepackage{fancyhdr}
\usepackage{lastpage}
\pagestyle{fancy}
\fancyhf{}
\renewcommand{\headrulewidth}{0pt}
\fancyfoot[L]{\sffamily\footnotesize\color{muted}Aakarsh Anand \sep Curriculum Vitae \sep Updated October 2026}
\fancyfoot[R]{\sffamily\footnotesize\color{muted}\thepage\ of \pageref{LastPage}}

% A little more air than the one-page resume
\titlespacing*{\section}{0pt}{8pt}{4.5pt}
\setlist*[itemize]{itemsep=2pt, topsep=2.5pt}
\newcommand{\subhead}[1]{{\sffamily\color{accent}#1}\par}

\begin{document}

\header

% ---------- education ----------
\section{Education}
\textbf{University of California, Los Angeles (UCLA)}\hfill\dates{Los Angeles, CA}\par
\vspace{1pt}
Ph.D., Computer Science \sep Advisor: Sriram Sankararaman\hfill\dates{2024 -- 2029 (expected)}\par
\note{\small\hspace{1em}Major field: Artificial Intelligence \sep Minor fields: Data Science Computing,
  Computational Systems Biology}\par
M.S., Computer Science \note{(earned during the Ph.D.)}\hfill\dates{2026}\par
B.S., Computer Science\hfill\dates{2020 -- 2024}\par

% ---------- research ----------
\section{Research Experience}
\role{Graduate Student Researcher}{Sankararaman Lab, UCLA \note{(undergraduate researcher 2022--24)}}{Jan 2022 -- present}

\thread{Foundation models for wearable sensors}{with Yuzhe Yang's lab, UCLA}
\begin{itemize}
  \item One of three co-first authors of \textbf{Inertia-1} (NeurIPS 2026). Pretrained motion foundation models
    from 5M to 100M parameters, using recurrent, Transformer, ViT, and wav2vec-style architectures, on
    \textbf{18.2M~hours} of accelerometry from \textbf{115K+~people} in NHANES and UK Biobank. Tested how sensor type,
    body placement, sampling rate, window length, architecture, model size, pretraining objective, and data scale
    affect performance on 15 downstream datasets covering activity recognition, freezing of gait, and disease
    prediction.
  \item Designed and ran the sensor-fusion study, which found that combining body placements and sensor types
    consistently beats single-sensor models.
  \item Co-first author of the follow-up journal paper (in prep.), which applies the model to diagnosis, prognosis,
    quantitative traits, and disease progression in five cohorts. Three of them (NHATS, ALS, and the Personalized
    Parkinson Project) are held out from pretraining. The paper also studies the genetics of the learned embeddings
    and early ways of pairing them with LLMs.
\end{itemize}

\thread{Deep learning for genetics}{}
\begin{itemize}
  \item Developing a training loss that makes a network's learned trait heritable and genetically correlated with
    disease, so a GWAS on it can find loci the original measurement misses. The genetic terms are computed exactly
    inside the training loop, and the loss works with any architecture.
\end{itemize}

\thread{Statistical genetics at biobank scale}{}
\begin{itemize}
  \item Leading a genome-wide test for gene--gene interactions in rare disease. It extends FAME to carriers of
    pathogenic variants to study how genetic background changes penetrance. Supported by an Anthropic AI for
    Science grant.
  \item Leading a single-cell variance-component model that separates cell-type-shared and cell-type-specific
    genetic and environmental effects on gene expression, and follows them along continuous cell states.
  \item Co-leading an atlas of genetic correlation across 3,570 UK Biobank trait pairs, partitioned by
    tissue-specific gene expression in 53 tissues.
  \item Second author of a gene--environment interaction study submitted to Nature Genetics: ran early
    simulations on which phenotype transformations reveal or hide G$\times$E, shaped the method with the first
    author, and ran real-data analyses.
  \item Ran simulations (calibration, power, runtime) for QuadKAST, FAME, and ENGINE and real-data analyses for
    FLEX, and added FAME's LD-block removal step.
\end{itemize}

\vspace{2pt}
\role{Research Intern (NSF REU)}{Bruins in Genomics Summer Program, UCLA}{Jun -- Sep 2022}
\begin{itemize}
  \item Used symbolic regression to explain gene--gene interactions learned by nonlinear models (first-author poster).
\end{itemize}

\vspace{2pt}
\role{Research Assistant}{Boutros Lab, UCLA}{Jan 2021 -- Apr 2022}
\begin{itemize}
  \item Added the HISAT2 aligner to the lab's Nextflow alignment pipeline, redesigned its configuration, and
    benchmarked it. The pipeline is now part of Metapipeline-DNA (Cell Reports Methods 2026).
\end{itemize}

% ---------- publications ----------
\sectionnote{* equal contribution}\section{Publications}
\interlinepenalty=10000 % never split a citation across pages
\subhead{Peer-reviewed}
\begin{itemize}
  \pub{Inertia-1: An Open Exploration of Wearable Motion Foundation Models}
      {Z.~Xu*, \me*, S.~Jiang*, C.~Zhuang, Z.~Shuai, S.~Sankararaman, Y.~Yang}
      {\vn{NeurIPS 2026}\sep\href{https://github.com/yang-ai-lab/Inertia-1}{code}}
      {https://arxiv.org/abs/2607.06617}
  \pub{A biobank-scale method for learning modulators of gene-environment interaction underlying human complex
      traits from multiple environmental exposures}
      {Z.~Liu, A.~Ramteke, \me, A.~Gorla, M.~Jeong, S.~Sankararaman}
      {\vn{Genome Research} 2026\sep selected from RECOMB 2026}
      {https://doi.org/10.1101/gr.282105.126}
  \pub{Metapipeline-DNA: A comprehensive germline and somatic genomics Nextflow pipeline}
      {Y.~Patel, C.~Zhu, T.~N.~Yamaguchi, N.~K.~Wang, \ldots, \me, \ldots, P.~C.~Boutros (43 authors)}
      {\vn{Cell Reports Methods} 2026}
      {https://doi.org/10.1016/j.crmeth.2026.101340}
  \pub{A biobank-scale test of marginal epistasis reveals genome-wide signals of polygenic interaction effects}
      {B.~Fu, A.~Pazokitoroudi, Z.~Shi, A.~Kar, A.~Xue, \me, P.~Anand, Z.~Liu, R.~Border, P.~Pajukanta,
      N.~Zaitlen, S.~Sankararaman}
      {\vn{Nature Genetics} 57(12):3175--3184, 2025}
      {https://doi.org/10.1038/s41588-025-02411-y}
  \pub{A scalable adaptive quadratic kernel method for interpretable epistasis analysis in complex traits}
      {B.~Fu*, P.~Anand*, \me*, J.~Mefford, S.~Sankararaman}
      {\vn{Genome Research} 34(9):1294--1303, 2024\sep RECOMB 2024}
      {https://doi.org/10.1101/gr.279140.124}
\end{itemize}
\vspace{3pt}
\subhead{Under review}
\begin{itemize}
  \pubnolink{Disentangling phenotype scale, amplification, and effect heterogeneity underlying polygenic GxE
      architecture of complex traits}
      {Z.~Liu, \me, X.~Qie, M.~Costantino, M.~Jeong, A.~Gorla, N.~Zaitlen, A.~Dahl, S.~Sankararaman}
      {submitted to \vn{Nature Genetics}}
  \pub{Comprehensive gene heritability estimation reveals the genetic architecture of rare coding variants
      underlying complex traits}
      {Z.~Liu, B.~Fu, M.~Jeong, P.~Anand, \me, S.~Jang, A.~Gorla, J.~Zhu, P.~Pajukanta, P.~F.~Palamara,
      N.~Zaitlen, R.~Border, S.~Sankararaman}
      {bioRxiv 2025\sep in revision at \vn{Nature Genetics}}
      {https://doi.org/10.1101/2025.10.07.681018}
\end{itemize}

% ---------- posters ----------
\section{Posters}
\begin{itemize}
  \pub{Explaining potential epistasis in genomic data using symbolic representations of complex black box models}
      {\me, P.~Anand}
      {Bruins in Genomics Summer Program, UCLA, August 2022}
      {https://aakarsh-anand.github.io/files/big_poster.pdf}
\end{itemize}

% ---------- awards ----------
\section{Awards, Fellowships, and Grants}
\role{Anthropic AI for Science Grant}{Rare Disease Research, \$50K in Claude credits}{2026}
\role{NIH T32 Predoctoral Fellowship}{UCLA Genomic Analysis Training Program}{2026 -- 2027}
\note{\small\hspace{1em}NHGRI T32 HG002536. One-year appointment, renewable.}\par
\role{Warren Alpert Computational Biology and AI Network Fellowship}{UCLA}{2024}
\note{\small\hspace{1em}Attended the Computational Genomics Summer Institute in 2024 and 2025.}\par
\role{NSF REU Scholarship}{Bruins in Genomics Summer Program, UCLA}{2022}
\role{Andy Grove Intel Scholarship}{Intel}{2020}
\role{Dean's Honors List}{UCLA}{2020 -- 2024}

% ---------- teaching ----------
\section{Teaching}
\role{Teaching Assistant}{UCLA}{2025 -- 2026}
\begin{itemize}
  \item Machine Learning Applications in Genetics (CS C124/224)\hfill\dates{Spring 2026}
  \item AI for Computational Genomics (CS CM121/221, $\sim$120 students)\hfill\dates{Spring 2025}\\
    When the first half of the course was rebuilt around AI, co-wrote three new homework assignments and the
    discussion slides with the other TAs.
  \item Machine Learning Applications in Biomedicine (DS BMED 205)\hfill\dates{Winter 2025}
\end{itemize}

% ---------- software ----------
\section{Software}
\begin{itemize}
  \item \href{https://github.com/yang-ai-lab/Inertia-1}{\textbf{Inertia-1}}: open-source wearable motion foundation
    models (co-developer).
  \item \href{https://github.com/sriramlab/FAME}{\textbf{FAME}}: added the LD-block removal step and support for
    missing phenotypes.
  \item \href{https://github.com/uclahs-cds/pipeline-align-DNA}{\textbf{pipeline-align-DNA}}: added the HISAT2
    aligner, plus fixes and container updates (8 merged pull requests).
\end{itemize}

% ---------- industry ----------
\section{Industry Experience}
\role{Software Engineering Intern}{LabCorp, Burlington, NC}{Jun -- Sep 2023}
\begin{itemize}
  \item Built reproducible Snakemake pipelines for high-volume genetic sequencing data.
\end{itemize}

% ---------- press ----------
\section{Press}
\begin{itemize}
  \item \href{https://www.uclahealth.org/news/release/researchers-develop-new-computational-tool-understand-how}{%
    ``Researchers develop a new computational tool to understand how genetic interactions impact human traits''}
    (on FAME). UCLA Health news release, December 2025.
\end{itemize}

% ---------- skills ----------
\section{Skills}
\textbf{Machine learning:} PyTorch, self-supervised pretraining, time-series models, large-scale evaluation\par
\textbf{Statistics and genetics:} variance-component models, method of moments, heritability and genetic
correlation, GWAS (plink2), LDSC, single-cell genetics\par
\textbf{Engineering:} Python, C++, SLURM/HPC, Snakemake, Nextflow, Docker, Git, Claude Code\par
\textbf{Data:} UK Biobank, NHANES, NHATS, OneK1K, GTEx\par

\end{document}
